% !TEX program = xelatex
% 系统开发工具基础 · 实验报告模板
% 用法：
%   1. 复制到 weekN/ 下，重命名为 weekN.tex
%   2. 修改下方「个人信息」一区的宏（周次、日期等）
%   3. 按「实验一」示例结构逐题填写；每题 = expbanner + reqbox + 过程 + resbox + trapbox
%   4. 编译两遍：xelatex weekN.tex && xelatex weekN.tex

\documentclass[11pt,a4paper,fontset=none]{ctexart}

\usepackage{amsmath}
\usepackage[margin=2.4cm]{geometry}
\usepackage[most]{tcolorbox}
\usepackage{booktabs}
\usepackage{caption}
\usepackage{enumitem}
\usepackage{listings}
\usepackage{xcolor}
\usepackage{fancyhdr}
\usepackage{fontspec}
\usepackage{fontawesome5}
\usepackage{tikz}
\usepackage{hyperref}
\usetikzlibrary{positioning}

% ---------- 个人信息（只需改这一区） ----------
\newcommand{\coursename}{系统开发工具基础}
\newcommand{\reportweek}{2}
\newcommand{\studentname}{彭俊皓}
\newcommand{\studentid}{25020007100}
\newcommand{\reportdate}{2026 年 X 月 X 日}

% ---------- 字体 ----------
\setCJKmainfont{Noto Serif CJK SC}
\setCJKsansfont{Noto Sans CJK SC}
\setCJKmonofont{Noto Sans CJK SC}
\setmainfont{Liberation Serif}
\setsansfont{Liberation Sans}
\setmonofont{Liberation Mono}

% ---------- 配色：四题四色（题数不足时其余颜色闲置） ----------
\definecolor{expone}{RGB}{41, 98, 168}   % 蓝
\definecolor{exptwo}{RGB}{46, 125, 50}   % 绿
\definecolor{expthree}{RGB}{214, 116, 30} % 橙
\definecolor{expfour}{RGB}{106, 76, 172}  % 紫

% ---------- 页眉页脚 ----------
\pagestyle{fancy}
\fancyhf{}
\fancyhead[L]{\small\color{gray}\coursename}
\fancyhead[R]{\small\color{gray}第 \reportweek 周实验报告}
\fancyfoot[C]{\thepage}
\renewcommand{\headrulewidth}{0.4pt}
\renewcommand{\headrule}{\color{gray!50}\hrule width\headwidth height 0.4pt}

% ---------- 自定义盒子 ----------
\newtcolorbox{reqbox}[1]{
  enhanced, breakable, boxrule=0.8pt, arc=2pt,
  left=8pt, right=8pt, top=5pt, bottom=5pt,
  colback=#1!6, colframe=#1!70!black,
  title={\small\bfseries\color{white}\faBook\ 题目要求}, coltitle=white,
  colbacktitle=#1!75!black
}
\newtcolorbox{resbox}[1]{
  enhanced, breakable, boxrule=0.8pt, arc=2pt,
  left=8pt, right=8pt, top=5pt, bottom=5pt,
  colback=#1!6, colframe=#1!70!black,
  title={\small\bfseries\color{white}\faCheckCircle\ 结果与验证}, coltitle=white,
  colbacktitle=#1!75!black
}
\newtcolorbox{trapbox}[1]{
  enhanced, breakable, boxrule=0.8pt, arc=2pt,
  left=8pt, right=8pt, top=5pt, bottom=5pt,
  colback=#1!6, colframe=#1!70!black,
  title={\small\bfseries\color{white}\faLightbulb\ 要点与注意事项}, coltitle=white,
  colbacktitle=#1!75!black
}

% 实验标题横幅
\newcommand{\expbanner}[3]{%
  \par\vspace{1.6em}
  \begin{tcolorbox}[enhanced, colback=#1!10, colframe=#1!85!black,
    boxrule=1pt, arc=3pt, left=8pt, right=8pt, top=6pt, bottom=6pt]
    {\Large\bfseries\color{#1!70!black}实验 #2\qquad #3}
  \end{tcolorbox}
  \par\vspace{0.2em}\noindent\ignorespaces
}
% 小节小标题
\newcommand{\sublabel}[2]{%
  \par\vspace{1em}\noindent{\bfseries\color{#1!75!black}#2}\par\nobreak\vspace{0.25em}\noindent\ignorespaces
}

% 总结横幅（无编号）
\newcommand{\concbanner}[2]{%
  \par\vspace{1.6em}
  \begin{tcolorbox}[enhanced, colback=#1!10, colframe=#1!85!black,
    boxrule=1pt, arc=3pt, left=8pt, right=8pt, top=6pt, bottom=6pt]
    {\Large\bfseries\color{#1!70!black}#2}
  \end{tcolorbox}
  \par\vspace{0.2em}\noindent\ignorespaces
}

% ---------- 代码/输出样式 ----------
\lstdefinestyle{shell}{
  language=bash, backgroundcolor=\color{gray!8},
  basicstyle=\ttfamily\footnotesize, keywordstyle=\color{expone!70!black},
  commentstyle=\color{green!50!black}, frame=single, rulecolor=\color{gray!45},
  breaklines=true, showstringspaces=false, columns=fullflexible, keepspaces=true,
  numbers=left, numberstyle=\tiny\color{gray}
}
\lstdefinestyle{output}{
  language={}, backgroundcolor=\color{gray!8},
  basicstyle=\ttfamily\footnotesize, frame=single, rulecolor=\color{gray!45},
  breaklines=true, showstringspaces=false, columns=fullflexible, keepspaces=true
}

\hypersetup{colorlinks=true, linkcolor=expfour!70!black, urlcolor=blue!60!black}
\captionsetup{font=small, labelfont=bf}

\begin{document}

% ================= 标题区 =================
\begin{center}
  {\LARGE\bfseries \coursename}\\[2pt]
  {\Large 第 \reportweek 周实验报告}\\[14pt]
  {\large 姓名：\studentname\qquad 学号：\studentid}\\[4pt]
  {\large \reportdate}\\[4pt]
  {\large 仓库：\url{https://github.com/pjh456/sys-devtool/}}\\[10pt]
  \begin{tcolorbox}[enhanced, colback=expone!6, colframe=expone!60!black,
    arc=2pt, boxrule=0.6pt, width=0.92\textwidth, center upper]
    \small 本报告记录第 \reportweek 周各道实验题的完成过程、关键命令与验证结果；
    全部结果均在本机重新执行验证。
  \end{tcolorbox}
\end{center}

\vspace{0.4em}
% 完成情况总览表：行数按实际题数增删
\begin{table}[htbp]
  \centering
  \caption{第 \reportweek 周实验完成情况}
  \label{tab:overview}
  \begin{tabular}{clll}
    \toprule
    实验 & 主题 & 关键工具 & 状态 \\
    \midrule
    1 & （填写主题） & （填写工具） & \textcolor{expone}{\faCheckCircle\ 完成} \\
    2 & （填写主题） & （填写工具） & \textcolor{exptwo}{\faCheckCircle\ 完成} \\
    3 & （填写主题） & （填写工具） & \textcolor{expthree}{\faCheckCircle\ 完成} \\
    4 & （填写主题） & （填写工具） & \textcolor{expfour}{\faCheckCircle\ 完成} \\
    \bottomrule
  \end{tabular}
\end{table}

% ================= 实验一（示例结构，照抄后替换内容） =================
\expbanner{expone}{一}{（填写实验主题）}

\begin{reqbox}{expone}
\begin{itemize}[leftmargin=1.4em, itemsep=1pt]
  \item 题目要求第一条；
  \item 题目要求第二条；
  \item 题目要求第三条。
\end{itemize}
\end{reqbox}

\sublabel{expone}{\faTerminal\ 实现步骤}
\begin{lstlisting}[style=shell]
# 关键命令
ls -la
\end{lstlisting}

\begin{resbox}{expone}
实际运行输出：
\begin{lstlisting}[style=output]
（粘贴终端输出）
\end{lstlisting}
对结果的验证说明。
\end{resbox}

\begin{trapbox}{expone}
\begin{itemize}[leftmargin=1.4em, itemsep=1pt]
  \item 易错点 / 关键技巧一；
  \item 易错点 / 关键技巧二。
\end{itemize}
\end{trapbox}

% ================= 实验二（复制上一节整段，改颜色为 exptwo，依此类推） =================
% \expbanner{exptwo}{二}{（填写实验主题）}
% ...
% \expbanner{expthree}{三}{（填写实验主题）}
% \expbanner{expfour}{四}{（填写实验主题）}

% ================= 总结 =================
\concbanner{expone}{总结与收获}
\begin{itemize}[leftmargin=1.4em, itemsep=2pt]
  \item （按题归纳本周收获，每条一个要点）；
  \item 补充说明（如方法对比、踩坑复盘）。
\end{itemize}

\end{document}
